%%%%%%%%%%%%%%%%%%%%%%%%%%%%%%%%%%%%%%%%%%%%%%%%%%%%%%%%%%%%%%%%%%%%%%
% Lớp 12 | Nguyên hàm, tích phân | Nguyên hàm
% Dạng 2D4?1-3: Nguyên hàm cơ bản hàm lượng giác
% Số câu: 67  (sinh tự động bởi tach_cau.py; từ nay thầy sửa trực tiếp file này)
%%%%%%%%%%%%%%%%%%%%%%%%%%%%%%%%%%%%%%%%%%%%%%%%%%%%%%%%%%%%%%%%%%%%%%

%%%%%%%%%%%%%%%% TRẮC NGHIỆM 4 PHƯƠNG ÁN %%%%%%%%%%%%%%%%

% nguồn: Phan1_Lop_12_Dai_so_Chuong_4
\begin{ex}%[2D4N1-3]%[ma:96520831]
	Hàm  số nào sau đây có một nguyên hàm là hàm số $ y=\sin 2x$?
	 \choice
		{$ y=\cos  2x$}
		{$ y=\dfrac{\cos  2x}{2}$}
		{\True $  y=2\cos 2x$}
		{$ y=\dfrac{-\cos 2x}{2}$} 
	\loigiai{ 
	Ta có  $\displaystyle\int{2\cos 2x}\mathrm{\,d}x=\sin 2x+C$.}
\end{ex}

% nguồn: Phan1_Lop_12_Dai_so_Chuong_4
\begin{ex}%[2D4N1-3]%[ma:635ce5d0]
	 Khẳng định nào sau đây đúng?
	\choice 
		{\True $\displaystyle\int(\cos x-2\sin x)\mathrm{\,d}x=\sin  x+2\cos x+C$}
		{$\displaystyle\int(\cos x-2\sin x)\mathrm{\,d}x=-\sin  x+2\cos x+C$}
		{$\displaystyle\int(\cos x-2\sin x)\mathrm{\,d}x=\sin x-2\cos  x+C$}
		{$\displaystyle\int(\cos  x-2\sin x)\mathrm{\,d}x=-\sin x-2\cos x+C$}
	\loigiai{ 
		$\displaystyle\int(\cos  x-2\sin x)\mathrm{\,d}x=\sin x+2\cos x+C$
	} 
\end{ex}

% nguồn: Phan1_Lop_12_Dai_so_Chuong_4
\begin{ex}%[2D4N1-3]%[ma:4a7b4cb4]
	Nguyên  hàm của hàm số $f(x)=\cos x$ là
	\choice 
		{$\cot  x+C$}
		{$\tan  x+C$}
		 {\True $\sin x+C$}
		{$-\sin  x+C$}
	\loigiai{ 
		Ta có $\displaystyle\int  f(x)\mathrm{\,d}x=\displaystyle\int \cos x\mathrm{\,d}x=\sin x+C$.
	 }
\end{ex}

% nguồn: Phan1_Lop_12_Dai_so_Chuong_4
\begin{ex}%[2D4N1-3]%[ma:7255f494]
	Hàm số nào sau đây là một nguyên hàm của  hàm số $ y=\sin 2x$?
	 \choice
		{$ y=\cos  2x$}
		{$  y=\dfrac{\cos 2x}{2}$}
		{$ y=2\cos 2x$} 
		{\True $ y=\dfrac{-\cos 2x}{2}$} 
	\loigiai{ 
	Ta  có $\displaystyle\int{\sin 2x}\mathrm{\,d}x=\dfrac{-\cos 2x}{2}+C$.}
\end{ex}

% nguồn: Phan1_Lop_12_Dai_so_Chuong_4
\begin{ex}%[2D4N1-3]%[ma:c3c7bf61]
	Với $C$ là hằng số, họ nguyên hàm của hàm số $f(x)=3\cos  x-1$ trên $\mathbb{R}$ là
	\choice 
		 {$3\sin x-1+C$}
		{$-3\sin  x-x+C$}
		{$3\sin  x+x+C$}
		{\True $3\sin x-x+C$} 
	 \loigiai{
		Ta có $\displaystyle\int f(x)\mathrm{\,d}x=\displaystyle\int (3\cos x-1)\mathrm{\,d}x=3\sin x -x +C$. 
	} 
\end{ex}

% nguồn: Phan1_Lop_12_Dai_so_Chuong_4
\begin{ex}%[2D4N1-3]%[ma:0490939b]
	Họ nguyên hàm của hàm  số $f(x)=\dfrac{1}{\cos ^2x}$ là
	\choice 
		 {$-\tan x+C$}
		{$\cot  x+C$}
		{$-\cot x+C$} 
		{\True  $\tan x+C$}
	\loigiai{ 
		Ta có  $\displaystyle\int \dfrac{1}{\cos^2x}\mathrm{\,d}x=\tan x+C$.
	} 
\end{ex}

% nguồn: Phan1_Lop_12_Dai_so_Chuong_4
\begin{ex}%[2D4N1-3]%[ma:6d8fad7a]
	Nếu hàm số $ y=\sin  x$ là một nguyên hàm của hàm số $ y=f(x)$ thì
	\choice 
		{$ f(x)=\sin x$} 
		{$  f(x)=-\cos x$}
		{$  f(x)=-\sin x$}
		{\True $  f(x)=\cos x$}
	\loigiai{ 
	Vì hàm số $ y=\sin x$ là một nguyên hàm của hàm số $ y=f(x)$  nên $(\sin x)'=f(x)$ hay $ f(x)=\cos x$.}
\end{ex}

% nguồn: Phan1_Lop_12_Dai_so_Chuong_4
\begin{ex}%[2D4N1-3]%[ma:3e3fc932]
	Hàm số nào sau đây là  một nguyên hàm của hàm số $f(x)=\dfrac{1}{\sin^2x}$?
	\choice 
		{$F(x)=\tan  x +C$}
		{$F(x)=-\tan x  +C$}
		{\True $F(x)=-\cot x  +C$}
		{$F(x)=\cot x +C$} 
	\loigiai{Theo công  thức nguyên hàm $\displaystyle\int \dfrac{1}{\sin^2x} \mathrm{\,d}x = -\cot x +C$.}
\end{ex}

% nguồn: Phan1_Lop_12_Dai_so_Chuong_4
\begin{ex}%[2D4N1-3]%[ma:7bbe07e8]
	Họ nguyên hàm của hàm số $y=\dfrac{1}{\cos ^2x}$ là 
	 \choice
		 {$\tan x$}
		{$-\cot x+C$} 
		{$-\cot x$} 
		{\True  $\tan x+C$}
	\loigiai{ 
		Họ nguyên hàm của hàm số $y=\dfrac{1}{\cos  ^2x}$ là $\tan x+C$.
	} 
\end{ex}

% nguồn: Phan1_Lop_12_Dai_so_Chuong_4
\begin{ex}%[2D4N1-3]%[ma:919ea291]
	Mệnh đề nào dưới  đây đúng?
	\choice 
		{\True $\displaystyle\int  \cos x\,\mathrm{d}x=\sin x+C$}
		 {$\displaystyle\int \cos x\,\mathrm{d}x=-\sin x+C$}
		{$\displaystyle\int  \cos x\,\mathrm{d}x=-\cos x+C$}
		 {$\displaystyle\int \cos x\,\mathrm{d}x=\frac{1}{2}\cos^2 x+C$}
	\loigiai{Ta có $\displaystyle\int \cos x\,\mathrm{d}x=\sin x+C$.} 
\end{ex}

% nguồn: Phan1_Lop_12_Dai_so_Chuong_4
\begin{ex}%[2D4N1-3]%[ma:975d4202]
	Cho $a, b \in \mathbb{R}, a \neq 0$. Hàm số $y = \sin(ax+b)$  là một nguyên hàm của hàm số
	 \choice
		{\True $y  = a \cos(ax+b)$}
		{$y = -\dfrac{\cos(ax+b)}{a}$} 
		{$y =  -a \cos(ax+b)$}
		{$y  = \dfrac{\cos(ax+b)}{a}$}
	\loigiai{ 
		 Ta có $y' = (\sin(ax+b))'=a \cos(ax+b)$.
	 }
\end{ex}

% nguồn: Phan1_Lop_12_Dai_so_Chuong_4
\begin{ex}%[2D4N1-3]%[ma:65e04ee1]
	Họ nguyên hàm của hàm số  $y=\sin x$ là
	\choice 
		{\True  $-\cos x+C$}
		{$-\sin x+C$} 
		{$\cos  x+C$}
		 {$\tan x+C$}
	\loigiai{ 
		Họ  nguyên hàm của hàm số $y=\sin x$ là $-\cos x+C$.
	 }
\end{ex}

% nguồn: Phan1_Lop_12_Dai_so_Chuong_4
\begin{ex}%[2D4N1-3]%[ma:84d73b98]
	Nguyên hàm  của hàm số $f(x)=x-\sin x$ là
	\choice 
		{$\dfrac{x^2}{2}-\cos x+C$} 
		{\True  $\dfrac{x^2}{2}+\cos x+C$}
		{$x^2-\cos  x+C$}
		{$x^2+\cos  x+C$}
	 \loigiai{
		Ta có $\displaystyle\int (x-\sin x)  \mathrm{\,d}x=\dfrac{x^2}{2}+\cos x+C$.
	} 
\end{ex}

% nguồn: Phan1_Lop_12_Dai_so_Chuong_4
\begin{ex}%[2D4N1-3]%[ma:6a8d6d55]
	 Nguyên hàm của hàm số $g(x)=\cos x$ là
	 \choice
		{$\displaystyle\int g(x)\mathrm{\,d}x=-\sin  x+C$}
		{$\displaystyle\int g(x)\mathrm{\,d}x=\cos  x+C$}
		 {$\displaystyle\int g(x)\mathrm{\,d}x=-\cos x+C$}
		{\True $\displaystyle\int g(x)\mathrm{\,d}x=\sin x+C$} 
	 \loigiai{
		Ta có $(\sin  x)'=\cos x=g(x)$ nên nguyên hàm của hàm số $g(x)=\cos x$ là $\displaystyle\int g(x)\mathrm{\,d}x=\sin x+C$.
	 }
\end{ex}

% nguồn: Phan1_Lop_12_Dai_so_Chuong_4
\begin{ex}%[2D4N1-3]%[ma:5026f45a]
	Cho $a, b \in \mathbb{R},  a \neq 0$. Hàm số $y = \cos(ax+b)$ là một nguyên hàm của hàm số
	\choice 
		{\True $y = -a  \sin(ax+b)$}
		{$y = -\dfrac{\sin(ax+b)}{a}$} 
		 {$y = a \sin(ax+b)$}
		{$y =  \dfrac{\sin(ax+b)}{a}$}
	\loigiai{ 
		 Ta có $y' = (\cos(ax+b))'=-a \sin(ax+b)$.
	} 
\end{ex}

% nguồn: Phan1_Lop_12_Dai_so_Chuong_4
\begin{ex}%[2D4N1-3]%[ma:0d269d51]
	Họ các nguyên hàm của hàm số $f(x)=\cos  x-\sin x$ là
	 \choice
		{$\cos x-\sin x+C$} 
		 {\True $\sin x+\cos x+C$}
		 {$\sin x-\cos x+C$}
		 {$\sin2x+C$}
	\loigiai{Ta có $\displaystyle\int f(x)\mathrm{\,d}x=\displaystyle\int(\cos x-\sin  x)\mathrm{\,d}x=\sin x+\cos x+C$.}
\end{ex}

% nguồn: Phan1_Lop_12_Dai_so_Chuong_4
\begin{ex}%[2D4H1-3]%[ma:f0214c99]
	 Biết $\displaystyle \int (2\sin x+\cos x)^2 \mathrm{\,d}x=a\sin 2x-\cos 2x+bx+C$, với $a, b\in \mathbb{Q}$. Tính $a^2+b^2$.
	 \choice
		{$\dfrac{17}{2}$} 
		 {$\dfrac{109}{4}$}
		 {$\dfrac{17}{16}$}
		{\True  $\dfrac{109}{16}$}
	\loigiai{ 
		Từ giả thiết, áp dụng định nghĩa nguyên hàm ta có 
		\[(a\sin 2x-\cos 2x+bx+C)'=(2\sin  x+\cos x)^2=2\sin 2x-\dfrac{3}{2}\cos 2x+\dfrac{5}{2}.\]
		Ta có \[(a\sin 2x-\cos  2x+bx+C)'=2a\cos 2x+2\sin 2x+b.\]
		Đối chiếu kết quả đồng nhất, ta được \[\heva{&2a=-\dfrac{3}{2}\\&b=\dfrac{5}{2}} 
		\Leftrightarrow  \heva{&a=-\dfrac{3}{4}\\&b=\dfrac{5}{2}.}\]
		Vậy \[a^2+b^2=\dfrac{9}{16}+\dfrac{25}{4}=\dfrac{109}{16}.\] 
	} 
\end{ex}

% nguồn: Phan1_Lop_12_Dai_so_Chuong_4
\begin{ex}%[2D4H1-3]%[ma:0589296b]
	Nguyên hàm của hàm  số $f(x)=2\cos^2\dfrac{x}{2}$ là
	 \choice
		{$4\cos  \dfrac{x}{2}+C$}
		{\True $x+\sin  x+C$}
		{$2\sin^2\dfrac{x}{2}+C$} 
		 {$\dfrac{2}{3}\cos^3\dfrac{x}{2}+C$}
	 \loigiai{
		Ta  có \[\displaystyle \int f(x)\mathrm{\,d}x=\int 2\cos^2\dfrac{x}{2}\mathrm{\,d}x
		=\int (1+\cos  x)\mathrm{\,d}x
		=\int \mathrm{\,d}x+\int  \cos x\mathrm{\,d}x
		=x+\sin x+C.\] 
	 }
\end{ex}

% nguồn: Phan1_Lop_12_Dai_so_Chuong_4
\begin{ex}%[2D4H1-3]%[ma:9ea0eb20]
	 Tìm $\displaystyle \int \sin 5x \mathrm{\,d}x$.
	 \choice
		{$\displaystyle \int \sin 5x \mathrm{\,d}x=-5\cos 5x+C$} 
		 {$\displaystyle \int \sin 5x \mathrm{\,d}x=-\cos 5x+C$}
		{\True $\displaystyle \int \sin  5x \mathrm{\,d}x=-\dfrac{1}{5}\cos 5x+C$}
		{$\displaystyle \int \sin 5x \mathrm{\,d}x=\dfrac{1}{5}\cos 5x+C$} 
	 \loigiai{
		Ta  có
		\[\displaystyle  \int \sin 5x ~\mathrm{\,d}x=-\dfrac{1}{5}\cos 5x+C.\]
	} 
\end{ex}

% nguồn: Phan1_Lop_12_Dai_so_Chuong_4
\begin{ex}%[2D4H1-3]%[ma:a3193d13]
	Hàm số nào sau đây là một nguyên  hàm của hàm số $ y=\tan ^2x$?
	 \choice
		 {$y=\dfrac{\tan^3x}{3}$}
		{\True $y=\tan  x-x$}
		 {$y=-\tan x+x$}
		 {$y=\tan x$}
	\loigiai{ 
	 Ta có $\displaystyle\int{\tan^2x}\mathrm{\,d}x=\displaystyle\int{\left[\left(\tan^2x+1\right)-1\right]}\mathrm{\,d}x=\displaystyle\int{\left(\dfrac{1}{\cos^2x}-1\right)}\mathrm{\,d}x=\tan x-x+C$.}
\end{ex}

% nguồn: Phan1_Lop_12_Dai_so_Chuong_4
\begin{ex}%[2D4H1-3]%[ma:e4717ac9]
	Một nguyên hàm của hàm số $f(x)  =\sin2x$ là
	 \choice
		{$F(x)=2\cos2x$} 
		{$F(x)=-\cos  2x$}
		 {$F(x)= \dfrac{1}{2}\cos 2x $}
		{$F(x)= \dfrac{-1}{2}\cos 2x $} 
	 \loigiai{\\
		Ta có: $F(x)= \displaystyle \int f(x)\mathrm{\,d}x = \displaystyle\int \sin{2x} \mathrm{\,d}x  = \dfrac{1}{2} \displaystyle \int \sin{2x} \mathrm{\,d}(2x) = -\dfrac{1}{2} \cos(2x) +C $ với $(C\in \mathbb{R})$.\\
		 Vậy ta chọn đáp án D.
	} 
\end{ex}

% nguồn: Phan1_Lop_12_Dai_so_Chuong_4
\begin{ex}%[2D4H1-3]%[ma:2b4c1a75]
	Tìm $\displaystyle \int \dfrac{\cos  2x}{\sin^2 x\cdot \cos^2 x} \mathrm{\,d}x$.
	 \choice
		 {$\cos x+\sin x+C$}
		{$-\cos x-\sin x+C$} 
		{\True  $-\cot x-\tan x+C$}
		 {$\cot x-\tan x+C$}
	 \loigiai{
		Ta có 
		\allowdisplaybreaks 
		\begin{eqnarray*}
			\displaystyle \int \dfrac{\cos 2x}{\sin^2 x\cdot \cos^2 x}  \mathrm{\,d}x
			&=&\int \dfrac{\cos^2  x-\sin^2 x}{\sin^2 x\cdot \cos^2 x} \mathrm{\,d}x\\
			&=&\int \left(\dfrac{1}{\sin^2 x}-\dfrac{1}{\cos^2 x}\right)\mathrm{\,d}x\\ 
			 &=&-\cot x-\tan x+C.
		\end{eqnarray*}
	 }
\end{ex}

% nguồn: Phan1_Lop_12_Dai_so_Chuong_4
\begin{ex}%[2D4H1-3]%[ma:5029c4eb]
	Nguyên  hàm của hàm số $y=\tan^2 x$ là
	 \choice
		{\True $\tan  x-x+C$}
		{$-\tan  x-x+C$}
		{$-\tan x+x+C$} 
		{$\tan x+x+C$} 
	\loigiai{ 
		 Ta có
		\[\displaystyle \int \tan^2 x  ~\mathrm{\,d}x
		=\int  \left(\dfrac{1}{\cos^2 x}-1\right)\mathrm{\,d}x=\tan x-x+C.\]
	} 
\end{ex}

% nguồn: Phan1_Lop_12_Dai_so_Chuong_4
\begin{ex}%[2D4H1-3]%[ma:edf52001]
	Một nguyên hàm $F(x)$ của hàm số $f(x)=\sin x+\dfrac{1}{\cos^2 x}$ thỏa mãn $F\left(\dfrac{\pi}{4}\right)=\dfrac{\sqrt{2}}{2}$  là
	\choice 
		{$-\cos x+\tan x-\sqrt{2}+1$} 
		{\True  $-\cos x+\tan x+\sqrt{2}-1$}
		{$\cos x+\tan  x+\sqrt{2}-1$}
		 {$-\cos x+\tan x+C$}
	\loigiai{ 
		 Ta có
		\[\displaystyle F(x)=\int \left(\sin x+\dfrac{1}{\cos^2 x}\right)\mathrm{\,d}x =-\cos x+\tan  x+C.\]
		 Vì $F\left(\dfrac{\pi}{4}\right)=\dfrac{\sqrt{2}}{2}$ nên $C=\sqrt{2}-1$.\\
		Vậy $F(x)=-\cos x+\tan  x+\sqrt{2}-1$.
	} 
\end{ex}

% nguồn: Phan1_Lop_12_Dai_so_Chuong_4
\begin{ex}%[2D4H1-3]%[ma:da5d969c]
	 Hàm số nào sau đây là một nguyên hàm của hàm số $ y=\cot ^2x$là
	\choice 
		 {$y=\dfrac{\cot^3x}{3}$}
		{$y=\cot x-x$} 
		 {$y=\cot x$}
		{\True  $y=-\cot x-x$}
	 \loigiai{
	 Ta có $\displaystyle\int{\cot^2x}\mathrm{\,d}x=\displaystyle\int{\left[\left(\cot^2x+1\right)-1\right]}\mathrm{\,d}x=\displaystyle\int{\left(\dfrac{1}{\sin^2x}-1\right)}\mathrm{\,d}x=-\cot x-x+C$.}
\end{ex}

% nguồn: Phan1_Lop_12_Dai_so_Chuong_4
\begin{ex}%[2D4H1-3]%[ma:6422861f]
	Nguyên hàm của hàm số $f(x)=2\sin x+\dfrac{3}{\sin^2x}$ là 
	\choice 
		{\True  $-2\cos x-3\cot x+C$}
		{$2\cos  x-3\tan x+C$}
		{$-2\cos x+3\cot x+C$} 
		{$2\cos  x-3\cot x+C$}
	\loigiai{ 
		Ta có 
		 \[\displaystyle \int f(x)\mathrm{\,d}x
		=\int  \left(2\sin x+\dfrac{3}{\sin^2x}\right)\mathrm{\,d}x
		=2\int  \sin x\mathrm{\,d}x+3\int{\dfrac{1}{\sin^2x}\mathrm{\,d}x}
		=-2\cos  x-3\cot x+C.\]
	} 
\end{ex}

% nguồn: Phan1_Lop_12_Dai_so_Chuong_4
\begin{ex}%[2D4H1-3]%[ma:1f1c49bc]
	Nguyên hàm $F(x)$ của hàm số  $f(x)=2\sin x-\cos x$ thỏa mãn $F\left(\dfrac{\pi}{3}\right)=-\dfrac{\sqrt{3}}{2}$ là
	 \choice
		{$F(x)=2\cos  x-\sin x-1$}
		{$F(x)=2\cos  x+\sin x-1-\sqrt{3}$}
		{\True $F(x)=-2\cos x-\sin  x+1$}
		 {$F(x)=-2\cos x-\sin x-1$}
	\loigiai{ 
		Ta có\\ 
		$\displaystyle F(x)=\int (2\sin x-\cos x)\mathrm{~d}x=-2\cos  x-\sin x+C$.\\
		$F\left(\dfrac{\pi}{3}\right)=-2\cos  \dfrac{\pi}{3}-\sin \dfrac{\pi}{3}+C =-\dfrac{\sqrt{3}}{2}\Rightarrow C=1$.\\
		Vậy  $F(x)=-2\cos x-\sin x+1$.
	 }
\end{ex}

% nguồn: Phan1_Lop_12_Dai_so_Chuong_4
\begin{ex}%[2D4H1-3]%[ma:4b428cd7]
	Hàm số nào dưới đây là một nguyên hàm  của hàm số $f(x)=x+\sin (2x+1)$?
	 \choice
		{$F(x)=\dfrac{1}{2}x^2-\cos (2x+1)$} 
		{$F(x)=\dfrac{1}{2}x^2-2\cos (2x+1)$} 
		{$F(x)=\dfrac{1}{2}x^2+\dfrac{1}{2}\cos  (2x+1)$}
		 {\True $F(x)=\dfrac{1}{2}x^2-\dfrac{1}{2}\cos (2x+1)$}
	 \loigiai{
		Ta  có
		\begin{itemize}
			\item $\left(\dfrac{1}{2}x^2-\cos(2x+1)\right)'=x+2\sin  (2x+1)$.
			\item $\left(\dfrac{1}{2}x^2-2\cos(2x+1)\right)'=x+4\sin  (2x+1)$.
			\item $\left(\dfrac{1}{2}x^2+\dfrac{1}{2}\cos(2x+1)\right)'=x-\sin  (2x+1)$.
			 \item $\left(\dfrac{1}{2}x^2-\dfrac{1}{2}\cos (2x+1)\right)'=x+\sin (2x+1)$.
		\end{itemize}
		Vậy $F(x)=\dfrac{1}{2}x^2-\dfrac{1}{2}\cos (2x+1)$ là một nguyên  hàm của hàm số $f(x)=x+\sin (2x+1)$.
	 }
\end{ex}

% nguồn: Phan1_Lop_12_Dai_so_Chuong_4
\begin{ex}%[2D4H1-3]%[ma:2f9ec564]
	Chiều cao của các bạn học sinh nữ lớp $12\mathrm{B}$ được ghi  lại ở bảng sau
	\begin{center}
		\begin{tabular}{|m{3.5cm}|m{2.5 cm}|m{2.5 cm}|m{2.5 cm}|m{2.5 cm}|}
			 \hline
			\centering\arraybackslash  Chiều cao $(\mathrm{cm})$&
			 \centering\arraybackslash$[155;160)$&
			\centering\arraybackslash$[160;165)$& 
			\centering\arraybackslash$[165;170)$& 
			\centering\arraybackslash$[170;175)$\\ 
			\hline 
			\centering\arraybackslash Số học  sinh&
			 \centering\arraybackslash$2$&
			\centering\arraybackslash$5$& 
			 \centering\arraybackslash$8$&
			\centering\arraybackslash$5$\\ 
			 \hline
		\end{tabular}
	\end{center}
	Khoảng  tứ phân vị của mẫu số liệu ghép nhóm trên là
	 \choice
		{$5$} 
		 {$6$}
		 {\True $7$}
		 {$8$}
	 \loigiai{
		Cỡ mẫu  $n=20$.\\
		Gọi $x_1,; x_2;\ldots;x_{20}$ là mẫu số liệu gốc gồm chiều cao của các bạn học sinh nữ được sắp xếp theo  thứ tự không giảm.\\
		 Tứ phân vị thứ nhất của mẫu số liệu gốc là $Q_1=\dfrac{1}{2}(x_{5}+x_{6})\in [160;165)$.\\
		Do đó  tứ phân vị thứ nhất của mẫu số liệu ghép nhóm là $$Q_1=160+\dfrac{\dfrac{20}{4}-2}{5}\cdot(165-160)=163\mathrm{\,(cm)}.$$
		Tứ  phân vị thứ ba của mẫu số liệu gốc là $Q_3=\dfrac{1}{2}(x_{15}+x_{16})$.\\
		Mà  $x_{15}\in [165;170)$ và $x_{16}\in [170;175)$.\\
		Do đó tứ phân vị thứ ba của mẫu số liệu ghép  nhóm là $Q_3=170$.\\
		Khoảng  tứ phân vị $\Delta Q=Q_3-Q_1=170-163=7\mathrm{\,(cm)}$.
	 }
\end{ex}

% nguồn: Phan1_Lop_12_Dai_so_Chuong_4
\begin{ex}%[2D4H1-3]%[ma:895e44de]
	Biết $F(x)$ là một nguyên hàm của của hàm số $f(x)=\sin x$ và đồ thị hàm số  $y=F(x)$ đi qua điểm $M(0;1)$. Tính $F\left(\dfrac{\pi}{2}\right)$.
	 \choice
		{\True  $F\left(\dfrac{\pi}{2}\right)=2$}
		 {$F\left(\dfrac{\pi}{2}\right)=-1$}
		 {$F\left(\dfrac{\pi}{2}\right)=0$}
		{$F\left(\dfrac{\pi}{2}\right)=1$} 
	 \loigiai{
		Theo bài ra $f(x)=\sin x$ nên $F(x)=-\cos x+C$,  với $C$ là hằng số tùy ý.\\
		Đồ thị  hàm số $y=F(x)$ đi qua điểm $M(0;1)$ nên
		$1=-\cos  0+C\Leftrightarrow C=2\Rightarrow F(x)=-\cos x+2$. \\
		Do đó  $F\left(\dfrac{\pi}{2}\right)=2$.
	} 
\end{ex}

% nguồn: Phan1_Lop_12_Dai_so_Chuong_4
\begin{ex}%[2D4H1-3]%[ma:fdd527a3]
	Cho hàm số $f(x)=3x^2+\sin  x$. Khẳng định nào sau đây đúng?
	 \choice
		{\True $\displaystyle\int f(x) \mathrm{d}x=x^3-\cos  x+C$}
		 {$\displaystyle\int f(x) \mathrm{d}x=x^3+\cos x+C$}
		 {$\displaystyle\int f(x) \mathrm{d}x=6x-\cos x+C$}
		{$\displaystyle\int f(x)  \mathrm{d}x=6x+\cos x+C$}
	\loigiai{ 
		Ta có $\displaystyle\int  f(x) \mathrm{d}x=\displaystyle\int (3x^2+\sin x) \mathrm{d}x=x^3-\cos x+C$.
	} 
\end{ex}

% nguồn: Phan1_Lop_12_Dai_so_Chuong_4
\begin{ex}%[2D4H1-3]%[ma:28087ffa]
	Nguyên hàm của  hàm số $f(x)=\left(\cos \dfrac{x}{2}+\sin \dfrac{x}{2}\right)^2$ là
	\choice 
		{$\sin  \dfrac{x}{2}-\cos \dfrac{x}{2}+C$}
		{\True $x-\cos x+C$} 
		{$\left(\sin  \dfrac{x}{2}-\cos \dfrac{x}{2}\right)^2+C$}
		{$\sin \dfrac{x}{2}+\cos \dfrac{x}{2}+C$} 
	 \loigiai{
		 Ta có
		\allowdisplaybreaks 
		\begin{eqnarray*}
			\displaystyle \int f(x)\mathrm{\,d}x 
			&=&\int  \left(\cos \dfrac{x}{2}+\sin \dfrac{x}{2}\right)^2\mathrm{\,d}x\\
			&=&\int\left(\cos^2\dfrac{x}{2}+\sin^2\dfrac{x}{2}+2\cos\dfrac{x}{2}\cdot\sin\dfrac{x}{2}\right)\mathrm{\,d}x\\ 
			&=&\int  (1+\sin x)\mathrm{\,d}x\\
			&=&\int  \mathrm{\,d}x+\int \sin x\mathrm{\,d}x\\
			&=&x-\cos x+C. 
		\end{eqnarray*}
	 }
\end{ex}

% nguồn: Phan1_Lop_12_Dai_so_Chuong_4
\begin{ex}%[2D4H1-3]%[ma:dde886b4]
	Một nguyên hàm $F(x)$ của hàm  số $f(x)=\sin x+2\cos x$ biết $F\left(\dfrac{\pi}{2}\right)=0$ là
	 \choice
		{$F(x)=-2\sin x-\cos  x+2$}
		{$F(x)=2\sin x-\cos  x+2$}
		{$F(x)=\sin x-2\cos  x-2$}
		 {\True $F(x)=2\sin x-\cos x-2$}
	\loigiai{ 
		Ta có 
		\[\displaystyle \int  (\sin x+2\cos x)\mathrm{\,d}x=-\cos x+2\sin x+C.\]
		 Do $F\left(\dfrac{\pi}{2}\right)=0$ nên $C=-2$ nên $F(x)=2\sin x-\cos x-2$.
	 }
\end{ex}

% nguồn: Phan1_Lop_12_Dai_so_Chuong_4
\begin{ex}%[2D4H1-3]%[ma:448c4393]
	Nguyên hàm của hàm số có  $f(x)=\tan^2x+\cot^2x$ là
	\choice 
		 {$2\tan x+2\cot x+C$}
		 {$\dfrac{1}{3}\tan^3x+\dfrac{1}{3}\cot^3x+C$}
		{$=\tan x+\cot  x-2x+C$}
		{\True  $\tan x-\cot x-2x+C$}
	\loigiai{ 
		Ta có 
		\begin{align*}
			 \displaystyle \int f(x)\mathrm{\,d}x
			 &=\int \tan^2x\mathrm{\,d}x+\int\cot^2x\mathrm{\,d}x\\
			 &=\int(1+\tan^2x-1)\mathrm{\,d}x+\int(1+\cot^2x-1)\mathrm{\,d}x\\
			 &=\int \left(\dfrac{1}{\cos^2x}-1\right)\mathrm{\,d}x +\int\left(\dfrac{1}{\sin^2x}-1\right)\mathrm{\,d}x\\
			&=\int\dfrac{1}{\cos^2x}\mathrm{\,d}x+\int\dfrac{1}{\sin^2x}\mathrm{\,d}x-2\int\mathrm{\,d}x\\ 
			&=\tan x-\cot  x-2x+C.
		\end{align*}
	 }
\end{ex}

% nguồn: Phan1_Lop_12_Dai_so_Chuong_4
\begin{ex}%[2D4H1-3]%[ma:c1272d4d]
	 Họ nguyên hàm của hàm số $ y=\sin ^2\dfrac{x}{2}$ là
	\choice 
		{$\dfrac{x+\sin  x}{2}+C$}
		 {\True $\dfrac{x-\sin x}{2}+C$}
		 {$\dfrac{x-\cos x}{2}+C$}
		{$\dfrac{x+\cos x}{2}+C$} 
	 \loigiai{
	Ta có $\displaystyle\int{\sin^2\dfrac{x}{2}}\mathrm{\,d}x=\displaystyle\int{\dfrac{1-\cos x}{2}}\mathrm{\,d}x=\dfrac{1}{2}\displaystyle\int{(1-\cos x)}\mathrm{\,d}x=\dfrac{1}{2}(x-\sin x)+C$.} 
\end{ex}

% nguồn: Phan1_Lop_12_Dai_so_Chuong_4
\begin{ex}%[2D4H1-3]%[ma:ef4842db]
	Nguyên hàm của hàm số $f(x)=(\tan x+\cot  x)^2$ là
	 \choice
		{$2x+\tan x-\cot x+C$} 
		{$2x+\tan x+\cot x+C$} 
		{\True $\tan  x-\cot x+C$}
		{$\tan x+\cot x+C$} 
	\loigiai{ 
		Ta có 
		\allowdisplaybreaks 
		\begin{eqnarray*}
			\displaystyle\int (\tan x+\cot x)^2 \mathrm{\,d}x&=&\displaystyle\int  \left(\tan^2 x+\cot^2 x+2\right) \mathrm{\,d}x\\
			 &=&\displaystyle\int \left(\dfrac{1}{\cos^2x}+\dfrac{1}{\sin^2x}\right) \mathrm{\,d}x\\
			&=&\tan x-\cot  x+C.
		\end{eqnarray*}
	} 
\end{ex}

% nguồn: Phan1_Lop_12_Dai_so_Chuong_4
\begin{ex}%[2D4H1-3]%[ma:a236759f]
	Cho hàm số $f(x)$ thỏa  mãn $f'(x)=3-5\cos x$ và $f(0)=5$. Mệnh đề nào dưới đây đúng?
	 \choice
		{$f(x)=3x-5\sin x-5$} 
		 {$f(x)=3x+5\sin x+5$}
		 {$f(x)=3x+5\sin x+2$}
		{\True $f(x)=3x-5\sin  x+5$}
	 \loigiai{
		Ta  có
		\[\displaystyle  f(x)=\int f'(x)\mathrm{\,d}x
		=\int  (3-5\cos x)\mathrm{\,d}x=3x-5\sin x+C.\]
		Ta lại  có
		\[f(0)=5\Rightarrow 3\cdot 0-5\sin 0+C=5\Leftrightarrow  C=5.\]
		Vậy $f(x)=3x-5\sin  x+5$.
	} 
\end{ex}

% nguồn: Phan1_Lop_12_Dai_so_Chuong_4
\begin{ex}%[2D4H1-3]%[ma:71a3cbe5]
	Cho hàm số $f(x)=\dfrac{1}{x}+\dfrac{1}{\cos^2x}$. Chọn khẳng định đúng. 
	\choice 
		{\True $\displaystyle \int f(x)\mathrm{\,d}x=\ln  |x|+\tan x+C$}
		{$\displaystyle  \int f(x)\mathrm{\,d}x=\ln x+\tan x+C$}
		{$\displaystyle \int f(x)\mathrm{\,d}x=\ln x+\tan  |x|+C$}
		{$\displaystyle \int f(x)\mathrm{\,d}x=\ln  |x|+\tan x$}
	\loigiai{ 
		Ta  có
		\[\displaystyle  \int f(x)\mathrm{\,d}x
		=\int\left(\dfrac{1}{x}+\dfrac{1}{\cos^2x}\right)\mathrm{\,d}x 
		=\ln |x|+\tan x+C.\] 
	} 
\end{ex}

% nguồn: Phan1_Lop_12_Dai_so_Chuong_4
\begin{ex}%[2D4H1-3]%[ma:f534578e]
	Nguyên  hàm của hàm số $f(x)=3\cos x-4\sin x$ là
	 \choice
		{$3\sin x-4\cos x$} 
		{$-3\sin x+4\cos x$} 
		 {\True $3\sin x+4\cos x+C$}
		{$-3\sin  x+4\cos x+C$}
	\loigiai{ 
		Ta có 
		\[\displaystyle  \int f(x)\mathrm{\,d}x
		=\int (3\cos x-4\sin x)\mathrm{\,d}x 
		=3\int \cos x\mathrm{\,d}x-4\int \sin x\mathrm{\,d}x=3\sin x+4\cos  x+C.\]
	} 
\end{ex}

% nguồn: Phan1_Lop_12_Dai_so_Chuong_4
\begin{ex}%[2D4H1-3]%[ma:95755e4d]
	 Họ nguyên hàm của hàm số $f(x)=3x^2+\sin x$ là
	 \choice
		 {\True $x^3-\cos x+C$}
		{$6x-\cos x+C$} 
		{$x^3+C$} 
		{$x^3+\sin  x+C$}
	\loigiai{ 
		 Ta có \[\displaystyle \int f(x)\mathrm{\,d}x
		=\int (3x^2+\sin  x)\mathrm{\,d}x=x^3-\cos x+C.\]
	 }
\end{ex}

% nguồn: Phan1_Lop_12_Dai_so_Chuong_4
\begin{ex}%[2D4H1-3]%[ma:22793cef]
	Với $C$ là hằng số, họ nguyên hàm của hàm số $f(x)=2\cos 2x$ là 
	 \choice
		 {\True $\sin 2x+C$}
		 {$2\sin 2x+C$}
		{$-2\sin  2x+C$}
		 {$-\sin 2x+C$}
	\loigiai{ 
		Ta có 
		 \[\displaystyle \int f(x)\mathrm{\,d}x
		=\int 2\cos 2x ~\mathrm{\,d}x=\sin  2x+C.\]
	 }
\end{ex}

% nguồn: Phan1_Lop_12_Dai_so_Chuong_4
\begin{ex}%[2D4H1-3]%[ma:35f09e92]
	Khẳng định nào  đây \textbf{sai}?
	 \choice
		{$\displaystyle \int \sin x  \mathrm{\,d}x=-\cos x+C$}
		{\True $\displaystyle \int \cos x ~\mathrm{\,d}x=-\sin  x+C$}
		{$\displaystyle  \int \cos x ~\mathrm{\,d}x=\sin x+C$}
		{$\displaystyle  \int \cos x ~\mathrm{\,d}x=-\tan x+C$}
	\loigiai{ 
		Ta có $\displaystyle \int  \cos x \mathrm{\,d}x=\sin x+C$.
	 }
\end{ex}


%%%%%%%%%%%%%%%% ĐÚNG/SAI %%%%%%%%%%%%%%%%

% nguồn: Phan1_Lop_12_Dai_so_Chuong_4
\begin{ex}%[2D4H1-3]%[ma:32ef72de]
	 Cho hàm số $f(x)=\tan^2x$ và $\displaystyle F(x)=\int f(x)\mathrm{\,d}x$. Khi đó
	 \choiceTF
		{\True $F'(x)=\tan^2x$} 
		 {$\tan x$ là một nguyên hàm của $f(x)$}
		{\True $\displaystyle F(x)=\int  \left(\dfrac{1}{\cos^2x}-1\right)\mathrm{\,d}x$}
		{Biết $F(0)=0\Rightarrow  F(\pi)=0$}
	\loigiai{ 
		\begin{itemchoice}
			\itemch  Vì $F'(x)=f(x)=\tan^2x$.
			 \itemch Ta có
			\[(\tan x)'=\left(\dfrac{\sin x}{\cos x}\right)' 
			 =\dfrac{\cos^2x+\sin^2x}{\cos^2x}=\tan^2x+1.\]
			 \itemch Do $\tan^2x=\dfrac{1}{\cos^2x}-1$ nên
			\[\displaystyle  F(x)=\int f(x)\mathrm{\,d}x
			=\int \tan^2x\mathrm{\,d}x 
			=\int  \left(\dfrac{1}{\cos^2x}-1\right)\mathrm{\,d}x.\]
			\itemch  Ta có
			\[\displaystyle F(x)=\int \left(\dfrac{1}{\cos^2x}-1\right)\mathrm{\,d}x 
			=\tan  x-x+C.\]
			 Khi đó
			\[F(0)=C=0\Rightarrow F(x)=\tan  x-x\Rightarrow F(\pi)=-\pi.\]
		\end{itemchoice}
	 }
\end{ex}

% nguồn: Phan1_Lop_12_Dai_so_Chuong_4
\begin{ex}%[2D4H1-3]%[ma:45b47619]
	Đặt $\displaystyle F(x)=-\int \sin x\mathrm{~d}x$ và  thỏa mãn $F(0)=1$. Khi đó
	 \choiceTF
		{\True $F(x)=\cos x$} 
		{Hàm số $F(x)$ là hàm  số lẻ trên $\mathbb{R}$}
		{$F(\alpha)=\dfrac{\sin 2\alpha}{\sin \alpha}$, $\forall \alpha  \neq k\pi$, $k\in \mathbb{Z}$}
		 {$F\left(\dfrac{\pi}{2^{2\,025}}\right)\cdot F\left(\dfrac{\pi}{2^{2\,024}}\right)\cdot F\left(\dfrac{\pi}{2^{2\,023}}\right)\cdots F\left(\dfrac{\pi}{2^3}\right)\cdot F\left(\dfrac{\pi}{2^2}\right)=\dfrac{1}{\sin \dfrac{\pi}{2^{2\,025}}}$}
	\loigiai{ 
		\begin{itemchoice}
			\itemch Ta có $\displaystyle F(x)=-\int  \sin x\mathrm{~d}x=\cos x+C$.\\
			Mà $F(0)=1$ nên ta  có $C=0 \Rightarrow F(x)=\cos x$.
			\itemch Ta thấy rằng, $F(-x)=\cos  (-x)=\cos x=F(x)$, $\forall x\in \mathbb{R}$ nên $F(x)$ là hàm chẵn trên $\mathbb{R}$.
			\itemch Ta thấy $\sin 2\alpha =2\sin \alpha \cdot \cos \alpha \Rightarrow  \cos \alpha =\dfrac{\sin 2\alpha}{2\sin \alpha}$
			hay $F(\alpha)=\cos \alpha  =\dfrac{\sin 2\alpha}{2\sin \alpha}$.
			 \itemch Ta có
			\allowdisplaybreaks 
			\begin{eqnarray*}
				&&F\left(\dfrac{\pi}{2^{2\,025}}\right)\cdot F\left(\dfrac{\pi}{2^{2\,024}}\right)\cdot  F\left(\dfrac{\pi}{2^{2\,023}}\right)\cdots F\left(\dfrac{\pi}{2^3}\right)\cdot F\left(\dfrac{\pi}{2^2}\right)\\
				&=&\dfrac{\sin \dfrac{\pi}{2^{2\,024}}}{2\sin \dfrac{\pi}{2^{2\,025}}}\cdot \dfrac{\sin \dfrac{\pi}{2^{2\,023}}}{2\sin \dfrac{\pi}{2^{2\,024}}}\cdot \dfrac{\sin \dfrac{\pi}{2^{2\,022}}}{2\sin \dfrac{\pi}{2^{2\,023}}}\cdots  \dfrac{\sin \dfrac{\pi}{2^2}}{2\sin \dfrac{\pi}{2^3}}\cdot \dfrac{\sin \dfrac{\pi}{2}}{2\sin \dfrac{\pi}{2^2}}\\
				&=&\dfrac{1}{2^{2\,024}\sin\dfrac{\pi}{2^{2\,025}}}. 
			\end{eqnarray*}
		\end{itemchoice}
	 }
\end{ex}

% nguồn: Phan1_Lop_12_Dai_so_Chuong_4
\begin{ex}%[2D4H1-3]%[ma:f6fe9124]
	Cho hàm số $f(x)=\sin x+\cos x$.  Khi đó
	 \choiceTF
		{\True Ta  có $\displaystyle F(x)=\int f(x)\mathrm{~d}x=\sin x-\cos x+C$, với $C$ là hằng số}
		{\True Biết rằng,  $F(0)=-1$. Khi đó, $F(x)=\sin x-\cos x$}
		{\True Hàm số $F(x)$ đồng biến trên  khoảng $\left(0; \dfrac{\pi}{2}\right)$}
		{Hàm số $F(x)$ đạt  giá trị nhỏ nhất là $-2$}
	 \loigiai{
		\begin{itemchoice}
			\itemch  Ta có
			\[\displaystyle  F(x)=\int f(x)\mathrm{~d}x
			=\int (\sin x+\cos x)\mathrm{~d}x=-\cos x+\sin  x+C.\]
			\itemch  Từ giả thiết, ta được
			$\sin  0-\cos 0+C=-1\Leftrightarrow C=0$ nên $F(x)=\sin x-\cos x$.
			\itemch Ta có $F'(x)=\sin x+\cos x>0$, $\forall x\in \left(0;  \dfrac{\pi}{2}\right)$ nên $F(x)$ đồng biến trên khoảng $\left(0; \dfrac{\pi}{2}\right)$.
			\itemch Ta thấy \[\heva{&&\sin x\ge -1\\&&-\cos x\ge  -1}
			 \Rightarrow F(x)=\sin x-\cos x\geq -2.\]
			 Trong đó có $F(x)=-2\Leftrightarrow \heva{&\sin x=-1 \\&\cos x=1} \Rightarrow \sin^2x+\cos^2x=2\neq 1$.\\
			Từ đó suy ra hàm số không đạt giá trị nhỏ nhất là $-2$. 
		\end{itemchoice}
	} 
\end{ex}

% nguồn: Phan1_Lop_12_Dai_so_Chuong_4
\begin{ex}%[2D4H1-3]%[ma:d7e66fed]
	 Xét tính đúng sai của các khẳng định sau:
	\choiceTF 
		{$\displaystyle  \int \sin x\mathrm{\,d}x=\cos x$}
		 {\True $\displaystyle \int \cos x\mathrm{\,d}x=\sin x+C$}
		 {\True $\tan x+1$ là một nguyên hàm của hàm số $y=\dfrac{1}{\cos^2x}$}
		{$\cot (x+1)$  là một nguyên hàm của hàm số $y=-\dfrac{1}{\sin^2x}$}
	\loigiai{ 
		\begin{itemchoice}
			 \itemch
			 Ta có $\displaystyle \int \sin x\mathrm{\,d}x=-\cos x+C$.
			\itemch 
			Ta có $\displaystyle  \int \cos x\mathrm{\,d}x=\sin x+C$.
			\itemch Do $(\tan x+1)'=\dfrac{1}{\cos^2x}$ nên $\tan x+1$  là một nguyên hàm của hàm số $y=\dfrac{1}{\cos^2x}$.
			\itemch Do $\big[\cot (x+1)\big]'=-\dfrac{1}{\sin^2(x+1)}$ nên $\cot  (x+1)$ không phải là một nguyên hàm của hàm số $y=-\dfrac{1}{\sin^2x}$.
		\end{itemchoice}
	} 
\end{ex}

% nguồn: Phan1_Lop_12_Dai_so_Chuong_4
\begin{ex}%[2D4H1-3]%[ma:ab00a46b]
	Đặt $\displaystyle F(x)=\int \dfrac{\mathrm{~d}x}{\sin^2x}$. Khi  đó
	 \choiceTF
		 {\True Ta có $F(x)=-\cot x+C$, với $C$ là hằng số}
		 {\True Biết rằng $F\left(\dfrac{\pi}{4}\right)=0$. Khi đó, ta có $F(x)=1-\cot x$}
		{Hàm số $F(x)$ là hàm số chẵn trên  tập xác định của $F(x)$}
		{Phương trình $F(x)=0$ có các nghiệm là $x=\dfrac{\pi}{4}+k2\pi$, với $k\in \mathbb{Z}$} 
	 \loigiai{
		\begin{itemchoice}
			 \itemch Ta có
			\[\displaystyle  F(x)=\int \dfrac{\mathrm{~d}x}{\sin^2x}
			 =-\cot x+C.\]
			\itemch Từ  giả thiết, ta được $C-\cot \dfrac{\pi}{4}=0\Leftrightarrow C=1$ nên $F(x)=1-\cot x$.
			\itemch Ta thấy $F\left(\dfrac{\pi}{4}\right)=0$, $F\left(\dfrac{-\pi}{4}\right)=2$ nên $F\left(\dfrac{-\pi}{4}\right)\ne \pm F\left(\dfrac{\pi}{4}\right)$ nên  hàm số đã cho không là hàm chẵn mà cũng không là hàm lẻ trên tập xác định của $F(x)$.
			\itemch Ta có $F(x)=0\Leftrightarrow 1-\cot x=0\Leftrightarrow \cot  x=1
			\Leftrightarrow x=\dfrac{\pi}{4}+k\pi$, $k\in \mathbb{Z}$. 
		\end{itemchoice}
	 }
\end{ex}

% nguồn: Phan1_Lop_12_Dai_so_Chuong_4
\begin{ex}%[2D4H1-3]%[ma:b9b945ad]
	Cho hàm số $f(x)=\sin x$. Khi  đó
	\choiceTF 
		{\True  Ta có $\displaystyle \int f(x)\mathrm{~d}x=-\cos x+C$, với $C$ là hằng số}
		{\True Gọi $F(x)$ là một nguyên hàm của hàm số $f(x)$ trên $\mathbb{R}$ thỏa mãn $F(\pi)=1$. Khi đó  ta có $F(0)=-1$}
		 {Ta có $\displaystyle \int F(x)\mathrm{~d}x=\sin x+C_1$, với $C_1$ là hằng số}
		{\True Phương trình $F(x)=f(x)$ có đúng $4$ nghiệm trên  đoạn $\left[0; 4\pi\right]$}
	\loigiai{ 
		\begin{itemchoice}
			 \itemch Ta có
			\[\displaystyle  \int f(x)\mathrm{~d}x
			=\int \sin x\mathrm{~d}x=-\cos  x+C.\]
			\itemch Từ giả thiết, ta  được
			\[-\cos \pi +C=1\Leftrightarrow 1+C=1\Leftrightarrow C=0.\] 
			Suy  ra $F(x)=-\cos x\Rightarrow F(0)=-1$.
			\itemch Ta  có
			\[\displaystyle \int  F(x)\mathrm{~d}x
			=-\int \cos x\mathrm{~d}x=-\sin x+C_1.\] 
			\itemch Ta có $F(x)=f(x)\Leftrightarrow  -\cos x=\sin x
			\Leftrightarrow \tan x=-1\Leftrightarrow x=\dfrac{-\pi}{4}+k\pi$,  $k\in \mathbb{Z}$.\\
			Ta  xét $0\le \dfrac{-\pi}{4}+k\pi \le 4\pi \Leftrightarrow \dfrac{1}{4}\le k\le \dfrac{17}{4}$.\\
			Vì $k\in  \mathbb{Z}$ nên $k\in \left\{1;2; 3; 4\right\}$.\\
			 Tức là phương trình đã cho có đúng $4$ nghiệm thuộc đoạn $\left[0;4\pi\right]$.
		\end{itemchoice}
	 }
\end{ex}

% nguồn: Phan1_Lop_12_Dai_so_Chuong_4
\begin{ex}%[2D4H1-3]%[ma:7723e926]
	Cho hàm số $f(x)=\cos^2x$ và $\displaystyle F(x)=\int{f(x)\mathrm{\,d}x}$. Khi đó 
	\choiceTF 
		 {$F'(x)=\sin 2x$}
		{\True $\displaystyle F(x)=\int  \dfrac{1+\cos 2x}{2}\mathrm{\,d}x$}
		{\True  $\dfrac{\sin x\cos x+x}{2}$ là một nguyên hàm của hàm số $f(x)$}
		{Biết $F(0)=0$.  Suy ra $F(\pi)=\pi$}
	 \loigiai{
		\begin{itemchoice}
			\itemch  Ta có $F'(x)=f(x)=\cos^2x$.
			 \itemch Do $f(x)=\cos^2x=\dfrac{1+\cos 2x}{2}$.
			\itemch Ta có \[\left(\dfrac{\sin  x\cos x+x}{2}\right)'
			 =\dfrac{\cos^2x-\sin^2x+1}{2}=\dfrac{2\cos^2x}{2}=\cos^2x.\]
			\itemch Do $\dfrac{\sin x\cos x+x}{2}$ là một nguyên hàm của hàm số $f(x)$ nên  $F(x)=\dfrac{\sin x\cos x+x}{2}+C$.\\
			Ta lại  có $F(0)=C=0\Rightarrow F(x)=\dfrac{\sin x\cos x+x}{2}
			\Rightarrow F(\pi)=\dfrac{\pi}{2}$. 
		\end{itemchoice}
	} 
\end{ex}

% nguồn: Phan1_Lop_12_Dai_so_Chuong_4
\begin{ex}%[2D4H1-3]%[ma:0db44aeb]
	Đặt $\displaystyle F(x)=\int \tan^2x\mathrm{~d}x$ và thỏa mãn  $F(0)=0$. Khi đó
	\choiceTF 
		 {\True Ta có $F(x)=-x+\tan x$}
		{Ta thấy  hàm số $F(x)$ là hàm số chẵn trên tập xác định của $F(x)$}
		{\True Hàm số $F(x)$ đồng biến  trên mỗi khoảng xác định của nó}
		{\True  Phương trình $F(x)+x=0$ có các nghiệm là $x=k\pi$, $k\in \mathbb{Z}$}
	 \loigiai{
		\begin{itemchoice}
			\itemch Ta có 
			 \[\displaystyle F(x)=\int \tan^2x\mathrm{~d}x
			=\int  \left(\dfrac{1}{\cos^2x}-1\right)\mathrm{~d}x
			=\tan x-x+C.\] 
			Mặt khác $F(0)=0$ nên ta  có $C=0$ nên $F(x)=-x+\tan x$.
			\itemch Ta thấy rằng $F(-x)=x-\tan x=-F(x)$ nên hàm số đã cho là hàm lẻ trên  mỗi khoảng xác định của $F(x)$.
			\itemch Ta có $F'(x)=\tan^2x\ge 0$ với mọi $x$  thuộc tập xác định của $F(x)$ nên hàm số $F(x)$ đồng biến trên mỗi khoảng xác định của nó.
			\itemch Ta  có
			\[F(x)+x=0\Leftrightarrow  -x+\tan x+x=0\Leftrightarrow \tan x=0\Leftrightarrow x=k\pi,\,k\in \mathbb{Z}.\]
		\end{itemchoice}
	} 
\end{ex}

% nguồn: Phan1_Lop_12_Dai_so_Chuong_4
\begin{ex}%[2D4H1-3]%[ma:1a73e98d]
	Cho hàm số $f(x)=\left(\sin{\dfrac{x}{2}}+\cos{\dfrac{x}{2}}\right)^2$. 
	 \choiceTF
		{\True  $f(x)=1+\sin{x}$}
		{\True $f(x)$ liên  tục trên $\mathbb{R}$}
		{$\displaystyle\int f(x)\mathrm{\,d}x=\displaystyle\int\mathrm{\,d}x+\displaystyle\int  (-\cos{x})\mathrm{\,d}x$}
		{$\displaystyle\int f(x)\mathrm{\,d}x=x+\cos{x}+C$} 
	 \loigiai{
		\begin{itemchoice}
			\itemch  \textbf{Đúng}.\\
			Ta  có
			\begin{align*}
				 f(x)=&\, \left(\sin{\dfrac{x}{2}}+\cos{\dfrac{x}{2}}\right)^2\\
				=&\,  \sin^2{\dfrac{x}{2}}+2\sin{\dfrac{x}{2}}\cos{\dfrac{x}{2}}+\cos^2{\dfrac{x}{2}}\\
				=&\, \sin^2{\dfrac{x}{2}}+\cos^2{\dfrac{x}{2}}+\sin{\left(2\cdot\dfrac{x}{2}\right)}\\ 
				 =&\, 1+\sin{x}.
			\end{align*}
			 \itemch \textbf{Đúng}.\\
			Hàm số $f(x)=1+\sin{x}$ là hàm số liên tục trên  $\mathbb{R}$.
			 \itemch \textbf{Sai}.\\
			Ta có $\displaystyle\int  f(x)\mathrm{\,d}x=\displaystyle\int (1+\sin{x})\mathrm{\,d}x =\displaystyle\int\mathrm{\,d}x+\displaystyle\int \sin{x}\mathrm{\,d}x$.
			 \itemch \textbf{Sai}\\
			Ta  có $\displaystyle\int f(x)\mathrm{\,d}x=\displaystyle\int (1+\sin{x})\mathrm{\,d}x =\displaystyle\int \mathrm{\,d}x+\displaystyle\int \sin{x}\mathrm{\,d}x= x-\cos{x}+C$.
		\end{itemchoice}
	} 
\end{ex}

% nguồn: Phan1_Lop_12_Dai_so_Chuong_4
\begin{ex}%[2D4V1-3]%[ma:ab4247db]
	Cho hàm số $f(x)=2\cos x$ và $g(x)=2\sin^2 \dfrac{x}{2}$. Khi  đó
	\choiceTF 
		{\True $\displaystyle\int f(x)\mathrm{\,d}x=2\sin  x+C$}
		{$\displaystyle\int g(x)\mathrm{\,d}x=-\cos x+C$} 
		{\True $\displaystyle\int \left[f(x)+g(x)\right]\mathrm{\,d}x=x+\sin x+C$} 
		 {$\displaystyle\int \dfrac{f(x)}{g(x)-1} \mathrm{\,d}x=2x+C$ (biết $x$ thoả mãn $g(x)\ne 1$)}
	 \loigiai{
		\begin{itemchoice}
			\itemch {\bf  Đúng}. \\
			$\displaystyle\int f(x)\mathrm{\,d}x=\displaystyle\int 2\cos x\mathrm{\,d}x=2\displaystyle\int \cos x\mathrm{\,d}x=2\sin x+C$. 
			\itemch {\bf Sai}.  \\
			$\displaystyle\int g(x)\mathrm{\,d}x=\displaystyle\int 2\sin^2 \dfrac{x}{2}  \mathrm{\,d}x=\displaystyle\int (1-\cos x)\mathrm{\,d}x=x-\sin x+C$.
			\itemch {\bf  Đúng}.\\
			$\displaystyle\int \left[f(x)+g(x)\right]\mathrm{\,d}x=\displaystyle\int f(x)\mathrm{\,d}x+\displaystyle\int g(x)\mathrm{\,d}x=x+\sin  x+C$.
			\itemch {\bf  Sai}. \\
			$\displaystyle\int \dfrac{f(x)}{g(x)-1} \mathrm{\,d}x=\displaystyle\int \dfrac{2\cos x}{2\sin^2 \frac{x}{2}-1}  \mathrm{\,d}x=\displaystyle\int \dfrac{2\cos x}{-\cos x} \mathrm{\,d}x=\displaystyle\int (-2)\mathrm{\,d}x=-2x+C$.
		\end{itemchoice}
	} 
\end{ex}

% nguồn: Phan1_Lop_12_Dai_so_Chuong_4
\begin{ex}%[2D4V1-3]%[ma:096a6f06]
	Biết rằng $I=\displaystyle\int  \dfrac{\sin x \cos^2 x}{1+\cos x} \mathrm{\,d} x=a \cos x+b \cos 2x-\ln(1+\cos x)+C\ (a; b \in \mathbb{R})$.
	 \choiceTF
		 {$\sin x \mathrm{\,d} x=\mathrm{\,d}(\cos x)$}
		{\True $I=-\displaystyle\int \dfrac{\cos^2  x}{1+\cos x} \mathrm{\,d}{(\cos x)}$}
		{$I=\displaystyle\int\left(\cos  x+1-\dfrac{1}{1+\cos x}\right) \mathrm{\,d}(\cos x)$}
		{\True Giá trị  của $a+b$ là $\dfrac{3}{4}$}
	 \loigiai{
		\begin{itemchoice}
			\itemch {\bf  Sai}.\\
			$\sin x  \mathrm{\,d} x=-\mathrm{\,d}(\cos x)$.
			\itemch {\bf Đúng}.\\ 
			Ta có $I=\displaystyle\int \dfrac{\sin x \cos^2 x}{1+\cos x} \mathrm{\,d} x=-\displaystyle\int \dfrac{\cos^2  x}{1+\cos x} \mathrm{\,d}{(\cos x)}$.
			\itemch  {\bf Sai}.\\
			$I=-\displaystyle\int \dfrac{\cos^2  x}{1+\cos x} \mathrm{\,d}{(\cos x)}=-\displaystyle\int \left(\cos x-1+\dfrac{1}{1+\cos x}\right) \mathrm{\,d}{(\cos x)}=\displaystyle\int\left(-\cos x+1-\dfrac{1}{1+\cos x}\right) \mathrm{\,d}(\cos x)$.
			\itemch {\bf Đúng}.\\ 
			\begin{eqnarray*}
				I&=&\displaystyle\int\left(-\cos x+1-\dfrac{1}{1+\cos  x}\right) \mathrm{\,d}(\cos x)\\
				&=&-\dfrac{\cos^2 x}{2}+\cos x-\ln(1+\cos  x)+C_1 \\
				 &=&-\dfrac{1}{4} \cos 2 x+\cos x-\ln(1+\cos x)+C_1-\dfrac{1}{4}\\
				&=&\cos x-\dfrac{1}{4}  \cos 2 x-\ln(1+\cos x)+C; C=C_1-\dfrac{1}{4}.
			\end{eqnarray*}
			Do đó  $a=1$, $b=\dfrac{-1}{4} \Rightarrow a+b=\dfrac{3}{4}$.
		\end{itemchoice}
	 }
\end{ex}


%%%%%%%%%%%%%%%% TRẢ LỜI NGẮN %%%%%%%%%%%%%%%%

% nguồn: Phan2_Lop_12_Dai_so_Chuong_4
\begin{bt}%[2D4H1-3]%[ma:0eda5322]
	Biết nguyên hàm $\displaystyle I=\int \dfrac{\mathrm{\,d}x}{2\tan x+1}$ bằng $\dfrac{a}{b}\ln\left| 2\sin x+\cos x\right|+\dfrac{c}{b}x+C$ với $\dfrac{a}{b}$, $\dfrac{c}{b}$ là các phân số tối  giản và $a\in \mathbb{Z}$, $c\in \mathbb{Z}$, $b\in \mathbb{N}^*$. Tính $a-b+c$.
	 \shortans{-2}
	\loigiai{ 
		Ta  có
		\allowdisplaybreaks 
		\begin{eqnarray*}
			 I=\displaystyle\int \dfrac{\mathrm{\,d}x}{2\tan x+1}
			&=&\displaystyle\int \dfrac{\cos  x}{2\sin x+\cos x}\mathrm{\,d}x\\
			&=&\dfrac{1}{2}\displaystyle\int \dfrac{2\cos x-\sin  x+\sin x}{2\sin x+\cos x}\mathrm{\,d}x\\
			&=&\dfrac{1}{2}\displaystyle\int  \dfrac{2\cos x-\sin x}{2\sin x+\cos x}\mathrm{\,d}x+\dfrac{1}{2}\displaystyle\int \dfrac{\sin x}{2\sin x+\cos x}\mathrm{\,d}x.\\
		\end{eqnarray*}
		Đặt $J=\displaystyle\int \dfrac{\sin x}{2\sin  x+\cos x}\mathrm{\,d}x$ ta có $I=\dfrac{1}{2}\ln\left| 2\sin x+\cos x\right|+\dfrac{1}{2}J$.\\
		Ta có $\displaystyle 2J+I=\int  1\mathrm{\,d}x=x+C$ suy ra $J=\dfrac{1}{2}\left(x-I+C_1\right)$.\\
		 Khi đó ta có
		\allowdisplaybreaks 
		\begin{eqnarray*}
			 &&I=\dfrac{1}{2}\ln\left| 2\sin x+\cos x\right|+\dfrac{1}{4}(x-I+C_1)\\
			&\Leftrightarrow& I=\dfrac{4}{5}\left[\dfrac{1}{2}\ln\left| 2\sin x+\cos x\right|+\dfrac{1}{4}x\right]+C\\ 
			&\Leftrightarrow& I=\dfrac{2}{5}\ln\left| 2\sin x+\cos x\right|+\dfrac{1}{5}x+C. 
		\end{eqnarray*}
		Suy ra $a=2$,  $b=5$, $c=1$.\\
		Vậy  $a-b+c=-2$.
	 }
\end{bt}

% nguồn: Phan2_Lop_12_Dai_so_Chuong_4
\begin{bt}%[2D4H1-3]%[ma:b9064a82]
	Cho hàm số $y=f(x)$ có đạo hàm là $f'(x)=\sin 2x$,  $\forall x\in \mathbb{R}$ và $f\left(\dfrac{\pi}{4}\right)=0$. Biết $F(x)$ là nguyên hàm của $f(x)$ thỏa mãn $F\left(\dfrac{\pi}{2}\right)=2$, khi đó $F\left(\dfrac{\pi}{4}\right)$ bằng bao nhiêu?
	 \shortans{1,75}
	\loigiai{ 
		Ta  có
		\[f'(x)=\sin 2x,\,\forall x\in \mathbb{R} 
		\Rightarrow f(x)=\displaystyle \int \sin  2x\mathrm{\,d}x
		=-\dfrac{1}{2}\cos 2x+C.\] 
		 Vì $f\left(\dfrac{\pi}{4}\right)=0
		 \Leftrightarrow -\dfrac{1}{2}\cos \left(2\cdot \dfrac{\pi}{4}\right)+C=0
		\Leftrightarrow  C=0$ nên $f(x)=-\dfrac{1}{2}\cos 2x$.\\
		Vì $F(x)$ là nguyên hàm của $f(x)$  nên
		 $\displaystyle F(x)=\int \left(-\dfrac{1}{2}\cos 2x\right)\mathrm{\,d}x
		=-\dfrac{1}{4}\sin 2x+C$. \\ 
		Lại có 
		\[F\left(\dfrac{\pi}{2}\right)=2 
		\Leftrightarrow -\dfrac{1}{4}\sin  \left(2\cdot \dfrac{\pi}{2}\right)+C=2
		 \Leftrightarrow C=2.\]
		Suy  ra $F(x)=-\dfrac{1}{4}\sin 2x+2$.\\
		Khi  đó
		\[F\left(\dfrac{\pi}{4}\right)=-\dfrac{1}{4}\sin \left(2\cdot \dfrac{\pi}{4}\right)+2=\dfrac{7}{4}=1{,}75.\] 
	} 
\end{bt}

% nguồn: Phan2_Lop_12_Dai_so_Chuong_4
\begin{bt}%[2D4H1-3]%[ma:3bfd3303]
	Cho hàm số $y=f(x)$ có đạo  hàm là $f'(x)=4\sin 2x+\cos x$, $\forall x\in \mathbb{R}$ và $f(0)=-2$. Biết $F(x)$ là nguyên hàm của $f(x)$ thỏa mãn $F(\pi)=3$, khi đó $F\left(\dfrac{\pi}{2}\right)$ bằng bao nhiêu?
	 \shortans{2}
	\loigiai{ 
		Ta có 
		\[\displaystyle f(x)=\int  f'(x)\mathrm{\,d}x
		=\int (4\sin 2x+\cos  x)\mathrm{\,d}x
		=-2\cos 2x+\sin  x+C.\]
		Với $f(0)=-2\Rightarrow -2\cdot \cos 2\cdot 0+\sin 0+C=-2\Rightarrow C=0$ nên $f(x)=-2\cos 2x+\sin x$.  \\
		 Ta có
		 \[\displaystyle F(x)=\int f(x)\mathrm{\,d}x
		=\int  (-2\cos 2x+\sin x)\mathrm{\,d}x=-\sin 2x-\cos x+C'.\]
		Với $F(\pi)=3\Rightarrow  -\sin 2\pi -\cos \pi +C'=3\Rightarrow C'=2$ nên $F(x)=-\sin 2x-\cos x+2$. \\
		 Khi đó
		\[F\left(\dfrac{\pi}{2}\right)=-\sin \pi -\cos  \dfrac{\pi}{2}+2=2.\]
	 }
\end{bt}

% nguồn: Phan2_Lop_12_Dai_so_Chuong_4
\begin{bt}%[2D4H1-3]%[ma:0c88b190]
	Cho $F(x)=\cos 2x$ là một nguyên hàm của hàm số $f(x)\cdot \tan x$. Biết nguyên hàm của hàm số $\dfrac{\tan  x}{f'(x)}$ có dạng $\dfrac{a}{b}\tan x+C$ với $a$, $b$ là các số nguyên dương, $\dfrac{a}{b}$ là phân số tối giản và $C\in \mathbb{R}$. Tính $a-2b$.
	 \shortans{-15}
	\loigiai{ 
		Ta có $F'(x)=-2\sin 2x=-4\sin x\cos x$.\\ 
		Vì $F(x)=\cos 2x$ là một nguyên hàm của hàm  số $f(x)\cdot \tan x$ nên
		\[F'(x)=f(x)\cdot \tan x\Leftrightarrow -4\sin x\cos x=f(x)\cdot  \tan x\Leftrightarrow f(x)=-4\cos^2x.\]
		Suy  ra $f'(x)=8\sin x\cos x$.\\
		Khi đó  ta có
		\allowdisplaybreaks 
		\begin{eqnarray*}
			 \displaystyle \int \dfrac{\tan x}{f'(x)}\mathrm{\,d}x
			&=&\int \dfrac{\tan  x}{8\sin x\cos x}\mathrm{\,d}x\\
			&=&\int \dfrac{1}{8\cos^2x}\mathrm{\,d}x\\ 
			&=&\dfrac{1}{8}\tan  x+C.
		\end{eqnarray*}
		Suy ra $a=1$,  $b=8$.\\
		Vậy $a-2b=-15$. 
	} 
\end{bt}

% nguồn: Phan2_Lop_12_Dai_so_Chuong_4
\begin{bt}%[2D4H1-3]%[ma:bfa5a3e0]
	Gọi $F(x)$ là họ các nguyên hàm của hàm số  $f(x)=8\sin 3x\cos x$. Biết rằng $F(x)$ có dạng $F(x)=a\cos 4x+b\cos 2x+C$. Khi đó $a-b$ bằng bao nhiêu?
	\shortans{1} 
	\loigiai{ 
		Ta  có
		\[f(x)=8\sin 3x\cos  x=4(\sin 4x+\sin 2x)=4\sin 4x+4\sin 2x.\]
		Khi đó 
		 \[F(x)=-\cos 4x-2\cos 2x+C \Rightarrow a=-1; b=-2.\]
		 Vậy $a-b=1$.
	} 
\end{bt}

% nguồn: Phan2_Lop_12_Dai_so_Chuong_4
\begin{bt}%[2D4H1-3]%[ma:1e24e4bc]
	Cho $F(x)$ là nguyên hàm của hàm số $f(x)=\sin^2x$  trên $\mathbb{R}$ thỏa mãn $F\left(\dfrac{\pi}{4}\right)=0$. Tính giá trị của biểu thức $S=F(-\pi)+2F\left(\dfrac{\pi}{2}\right)$ (kết quả làm tròn đến hàng phần mười).
	\shortans{-0,4} 
	\loigiai{ 
		Vì Cho $F(x)$ là nguyên hàm  của hàm số $f(x)=\sin^2x$ nên ta có\\
		\[\displaystyle \int \sin^2x\mathrm{\,d}x=\int\dfrac{1-\cos 2x}{2}\mathrm{\,d}x 
		 =\dfrac{1}{2}x-\dfrac{1}{4}\sin 2x+C=F(x).\]
		Ta  có
		\[F\left(\dfrac{\pi}{4}\right)=0 
		\Rightarrow  \dfrac{\pi}{8}-\dfrac{1}{4}+C=0
		\Rightarrow  C=\dfrac{1}{4}-\dfrac{\pi}{8}.\]
		 Suy ra
		 \[F(x)=\dfrac{1}{2}x-\dfrac{1}{4}\sin 2x+\dfrac{1}{4}-\dfrac{\pi}{8}.\]
		Khi  đó
		 \[S=F(-\pi)+2F\left(\dfrac{\pi}{2}\right)
		=\left(\dfrac{1}{4}-\dfrac{5\pi}{8}\right)+2\left(\dfrac{1}{4}+\dfrac{\pi}{8}\right) 
		=\dfrac{3}{4}-\dfrac{3\pi}{8} \approx -0{,}4.\] 
	} 
\end{bt}


%%%%%%%%%%%%%%%% TỰ LUẬN %%%%%%%%%%%%%%%%

% nguồn: Phan2_Lop_12_Dai_so_Chuong_4
\begin{bt}%[2D4H1-3]%[ma:f8dd8d10]
	 Tìm họ nguyên hàm sau:
	\begin{multicols}{3}
		\begin{enumerate}
			\item $\displaystyle \int  (2\sin x-3\cos x)\mathrm{\,d}x$;
			\item  $\displaystyle \int \cos^2\dfrac{x}{2}\mathrm{\,d}x$;
			\item $\displaystyle  \int \dfrac{\mathrm{\,d}x}{\sin^2x\cos^2x}$;
			 \item $\displaystyle \int (\sin x+\cos x)^2\mathrm{\,d}x$;
			\item $\displaystyle \int \dfrac{1}{\sin^2x}\mathrm{\,d}x$  trên $(0;\pi)$;
			\item  $\displaystyle \int 2\sin \dfrac{x}{2}\cos \dfrac{x}{2}\mathrm{\,d}x$.
		\end{enumerate}
	\end{multicols}
	 \loigiai{
		\begin{enumerate}
			\item Ta có 
			\[\displaystyle \int (2\sin x-3\cos x)\mathrm{\,d}x 
			=2\int \sin x\mathrm{\,d}x-3\int \cos  x\mathrm{\,d}x
			=-2\cos x-3\sin  x+C.\]
			 \item Ta có
			 \[\displaystyle \int \cos^2\dfrac{x}{2}\mathrm{\,d}x
			=\int \left(\dfrac{1+\cos  x}{2}\right)\mathrm{\,d}x
			=\dfrac{1}{2}\int  (1+\cos x)\mathrm{\,d}x
			 =\dfrac{1}{2}(x+\sin x)+C.\]
			\item Ta  có
			\[\displaystyle \int  \dfrac{\mathrm{\,d}x}{\sin^2x\cos^2x}
			 =\int \dfrac{\sin^2x+\cos^2x}{\sin^2x\cos^2x}\mathrm{\,d}x
			=\int\dfrac{\mathrm{\,d}x}{\cos^2x}+\int\dfrac{\mathrm{\,d}x}{\sin^2x}=\tan  x-\cot x+C.\]
			\item Ta có 
			 \[\displaystyle \int (\sin x+\cos x)^2\mathrm{\,d}x
			=\int\left(\sin^2x+2\sin  x\cos x+\cos^2x\right)\mathrm{\,d}x
			=\int (1+\sin 2x)\mathrm{\,d}x =x-\dfrac{1}{2}\cos 2x+C.\] 
			\item Vì $(-\cot  x)'=\dfrac{1}{\sin^2x}$, với mọi $x$ thuộc $(0;\pi)$ nên $F(x)=-\cot x$ là một nguyên hàm của $\dfrac{1}{\sin^2x}$ trên $(0;\pi)$ nên $\displaystyle \int \dfrac{1}{\sin^2x}\mathrm{\,d}x=-\cot x+C$ trên $(0;\pi)$.
			\item Ta  có
			 \[\displaystyle \int 2\sin \dfrac{x}{2}\cos\dfrac{x}{2}\mathrm{\,d}x=\int \sin x\mathrm{\,d}x=-\cos x+C.\]
		\end{enumerate}
	 }
\end{bt}

% nguồn: Phan2_Lop_12_Dai_so_Chuong_4
\begin{bt}%[2D4H1-3]%[ma:3902b965]
	Tìm: 
	\begin{listEX}[2]
		 \item[a)] $\displaystyle\int [4(2-3x)^2-3\cos x]\mathrm{\,d}x$;
		\item[b)] $\displaystyle\int \left(3x^3-\dfrac{1}{2x^3}\right)\mathrm{\,d}x$; 
		\item[c)] $\displaystyle\int  \left(\dfrac{2}{\sin^2 x}-\dfrac{1}{3\cos^2x}\right)\mathrm{\,d}x$;
		\item [d)]  $\displaystyle\int \left(3^{2x-2}+4\cos x\right)\mathrm{\,d}x$;
		\item [e)] $\displaystyle\int \left(4\sqrt[5]{x^4}+\dfrac{3}{\sqrt[3]{x^2}}\right)\mathrm{\,d}x$; 
		\item  [f)] $\displaystyle\int \left(\sin\dfrac{x}{2}-\cos\dfrac{x}{2}\right)^2\mathrm{\,d}x$.
	\end{listEX}
	\loigiai{ 
		\begin{enumerate}
			\item $\displaystyle\int  [4(2-3x)^2-3\cos x]\mathrm{\,d}x=\displaystyle\int [4(2-3x)^2]\mathrm{\,d}x-\displaystyle\int [3\cos x]\mathrm{\,d}x
			=4\displaystyle\int  (2-3x)^2\mathrm{\,d}x-3\displaystyle\int \cos x\mathrm{\,d}x
			=4\displaystyle\int (4-12x+9x^2)\mathrm{\,d}x-3\displaystyle\int \cos  x\mathrm{\,d}x
			=4\left(4x-6x^2+3x^3\right)-3\sin x+C=12x^3-26x^2+16x-3\sin  x+C$;
			 \item $\displaystyle\int \left(3x^3-\dfrac{1}{2x^3}\right)\mathrm{\,d}x=\displaystyle\int 3x^3\mathrm{\,d}x-\displaystyle\int \left(\dfrac{1}{2x^3}\right)\mathrm{\,d}x=3\displaystyle\int x^3\mathrm{\,d}x-\displaystyle\int \left(\dfrac{1}{2}x^{-3}\right)\mathrm{\,d}x=\dfrac{3}{4}x^4+\dfrac{1}{4}x^{-2}+C$;
			\item $\displaystyle\int  \left(\dfrac{2}{\sin^2 x}-\dfrac{1}{3\cos^2x}\right)\mathrm{\,d}x=-2\cot x-\dfrac{1}{3}\tan x+C$;
			\item $\displaystyle\int \left(3^{2x-2}+4\cos x\right)\mathrm{\,d}x=\dfrac{3^{2x-2}}{2\ln3}+4\sin  x+C$;
			 \item $\displaystyle\int \left(4\sqrt[5]{x^4}+\dfrac{3}{\sqrt[3]{x^2}}\right)\mathrm{\,d}x=\displaystyle\int \left(4x^{\frac{4}{5}}+3x^{-\frac{2}{3}}\right)\mathrm{\,d}x=\dfrac{20}{9}x^{\frac{9}{5}}+9x^{\frac{1}{3}}+C$;
			\item $\displaystyle\int \left(\sin\dfrac{x}{2}-\cos\dfrac{x}{2}\right)^2\mathrm{\,d}x=\displaystyle\int \left(\sin^2\dfrac{x}{2}-2\sin\dfrac{x}{2}\cos\dfrac{x}{2}+\cos^2\dfrac{x}{2}\right)\mathrm{\,d}x=\displaystyle\int \left(1-\sin x\right)\mathrm{\,d}x=x+\cos  x+C$.
		\end{enumerate}
	 }
\end{bt}

% nguồn: Phan2_Lop_12_Dai_so_Chuong_4
\begin{bt}%[2D4H1-3]%[ma:aef6a77d]
	Tìm họ tất cả  các nguyên hàm của hàm số sau:
	\begin{enumerate}
		\item $y=  2^x - \dfrac{1}{x}$;
		\item $y= x\sqrt{x} + 3\cos  x - \dfrac{2}{\sin^2 x}$.
	\end{enumerate}
\end{bt}

% nguồn: Phan2_Lop_12_Dai_so_Chuong_4
\begin{bt}%[2D4H1-3]%[ma:f038bf7a]
	 Tìm họ nguyên hàm của hàm số $f(x)$, biết:
	\begin{multicols}{3}
		\begin{enumerate}
			\item $f(x)=4x+3\cos  x$;
			\item $f(x)=2x+\dfrac{1}{\sin^2  x}$;
			\item  $f(x)=\left(\sin \dfrac{x}{2}+\cos \dfrac{x}{2}\right)^2$;
			\item  $f(x)=\sin 3x$;
			\item $f(x)=\cos  (2x+3)$;
			 \item $f(x)=\cos 5x\cdot \cos x$.
		\end{enumerate}
	\end{multicols}
	\loigiai{ 
		\begin{enumerate}
			\item Ta có $\displaystyle \int  f(x)\mathrm{\,d}x
			=\int  (4x+3\cos x)\mathrm{\,d}x =\dfrac{4x^2}{2}+3\sin x+C=2x^2+3\sin x+C$.
			\item Ta có $\displaystyle \int \left(2x+\dfrac{1}{\sin^2  x}\right)\mathrm{\,d}x =x^2-\cot x+C$.
			\item Ta có $\left(\sin  \dfrac{x}{2}+\cos \dfrac{x}{2}\right)^2
			=\sin^2 \dfrac{x}{2}+2\sin \dfrac{x}{2}\cos  \dfrac{x}{2}+\cos^2 \dfrac{x}{2}=1+\sin x$. \\
			Do đó $\displaystyle \int \left(\sin \dfrac{x}{2}+\cos \dfrac{x}{2}\right)^2 ~\mathrm{\,d}x =\int 1 \mathrm{\,d}x+\int \sin x  \mathrm{\,d}x=x-\cos x+C$.
			\item Ta có $\displaystyle  \int f(x)\mathrm{\,d}x
			=\int \sin 3x  \mathrm{\,d}x=-\dfrac{1}{3}\cos 3x+C$.
			 \item Ta có $\displaystyle \int f(x)\mathrm{\,d}x=\dfrac{1}{2}\sin (2x+3)+C$.
			\item Ta có $\displaystyle \int \cos 5x\cos x  ~\mathrm{\,d}x
			 =\dfrac{1}{2}\int (\cos 4x+\cos 6x)\mathrm{\,d}x
			=\dfrac{1}{2}\left(\dfrac{1}{4}\sin  4x+\dfrac{1}{6}\sin 6x\right)+C$.
		\end{enumerate}
	 }
\end{bt}

% nguồn: Phan2_Lop_12_Dai_so_Chuong_4
\begin{bt}%[2D4H1-3]%[ma:6d14d693]
	Cho hàm số $f(x)$ thỏa mãn $f'(x)=2-5\sin x$ và $f(0)=10$. Tìm hàm $f(x)$. 
	 \loigiai{
		 Ta có
		\[\displaystyle f(x)=\int (2-5\sin x)\mathrm{\,d}x=2x+5\cos  x+C.\]
		Mà $f(0)=10\Rightarrow  C=5$ nên $f(x)=2x+5\cos x+5$.
	 }
\end{bt}

% nguồn: Phan2_Lop_12_Dai_so_Chuong_4
\begin{bt}%[2D4H1-3]%[ma:3e9e42e6]
	Cho $F(x)=\cos 2x$ là một  nguyên hàm của hàm số $f(x)=\tan x$. Tìm họ nguyên hàm của hàm số $g(x)=\dfrac{\tan x}{f'(x)}$.
	 \loigiai{
		Ta  có \[F'(x)=-2\sin 2x=-4\sin x\cos x.\]
		Vì $F(x)=\cos 2x$ là một nguyên hàm của hàm số $f(x)\tan  x$ nên
		 \[F'(x)=f(x)\cdot \tan x
		\Leftrightarrow -4\sin  x\cos x=f(x)\cdot \tan x
		\Leftrightarrow  f(x)=-4\cos^2x.\]
		Suy ra $f'(x)=8\sin  x\cos x$.\\
		Vậy \[\displaystyle \int g(x)\mathrm{\,d}x 
		=\int  \dfrac{\tan x}{f'(x)}\mathrm{\,d}x
		=\int \dfrac{\tan  x}{8\sin x\cos x}\mathrm{\,d}x
		=\int \dfrac{1}{8\cos^2x}\mathrm{\,d}x=\dfrac{1}{8}\tan x+C.\] 
	} 
\end{bt}

% nguồn: Phan2_Lop_12_Dai_so_Chuong_4
\begin{bt}%[2D4V1-3]%[ma:040475c8]
	Một vật chuyển động với vận tốc là $v(t)=\dfrac{1}{2\pi}+\dfrac{\sin (\pi t)}{\pi}\, (\mathrm{m}/\mathrm{s})$. Gọi $S_1$ là quãng đường vật đó đi trong $2$ giây đầu và $S_2$ là quãng  đường đi từ giây thứ $3$ đến giây thứ $5$. Chứng minh $S_2>S_1$?
	\loigiai{ 
		Gọi  thời điểm ban đầu $S(0)=0;t=0;S=0\Rightarrow S_1=S(2)$. \\
		Khi  đó
		\[\displaystyle S(t)=\int \left[\dfrac{1}{2\pi}+\dfrac{\sin (\pi  t)}{\pi}\right]\mathrm{\,d}t
		 =\dfrac{1}{2\pi}t-\dfrac{1}{\pi^2}\cos (\pi t)+C.\]
		Ta có $S(0)=0\Rightarrow C-\dfrac{1}{\pi^2}=0\Rightarrow  C=\dfrac{1}{\pi^2}$. \\
		Suy ra  $S(t)=\dfrac{1}{2\pi}-\dfrac{1}{\pi^2}\cos (\pi t)+\dfrac{1}{\pi^2}$. \\
		 $S_1=S(2)=\dfrac{1}{2\pi}-\dfrac{1}{\pi^2}+\dfrac{1}{\pi^2}
		=\dfrac{1}{2\pi}\approx 0{,}159$;  \\
		$\displaystyle S_2=\int_3^5 v(t)\mathrm{\,d}t\approx  0{,}31
		\Rightarrow  S_2>S_1$ (đpcm).
	 }
\end{bt}

% nguồn: Phan2_Lop_12_Dai_so_Chuong_4
\begin{bt}%[2D4V1-3]%[ma:341a8e66]
	Một vật chuyển động với gia tốc phụ thuộc vào thời gian theo công thức $a(t)=\sin \left(2t+\dfrac{\pi}{3}\right)$. Biết tại thời điểm $t=0$ thì vận tốc và quãng đường đi được của vật đều bằng $0$. Thiết lập công thức tính quãng đường đi được của vật đó  theo thời gian.
	\loigiai{ 
		Vận tốc của  vật tại thời điểm $t$ là
		\[\displaystyle  v(t)=\int a(t)\mathrm{\,d}t
		=\int\sin \left(2t+\dfrac{\pi}{3}\right)\mathrm{\,d}t 
		=-\dfrac{1}{2}\cos  \left(2t+\dfrac{\pi}{3}\right)+C.\]
		Ta  có\\
		$v(0)=0\Leftrightarrow C=\dfrac{1}{4} 
		\Rightarrow v(t)=-\dfrac{1}{2}\cos \left(2t+\dfrac{\pi}{3}\right)  +\dfrac{1}{4}$.\\
		$\Rightarrow  s(t)=\displaystyle \int \left[-\dfrac{1}{2}\cos \left(2t+\dfrac{\pi}{3}\right)+\dfrac{1}{4}\right]\mathrm{\,d}t
		=-\dfrac{1}{4}\sin  \left(2t+\dfrac{\pi}{3}\right)+\dfrac{1}{4}t+D$.\\
		 Khi đó $s(0)=0\Leftrightarrow D=\dfrac{\sqrt{3}}{8}$ nên ta suy ra $s(t)=-\dfrac{1}{4}\sin \left(2t+\dfrac{\pi}{3}\right)+\dfrac{1}{4}t+\dfrac{\sqrt{3}}{8}$.
	} 
\end{bt}
